\section{Open Problems}
\label{sec:open}

Our result closes one question and sharpens several others.

\paragraph{What can a scheduler do that ordering cannot?}
\Cref{sec:arrival} shows that first-come-first-served admission captures 87.4\% of the reachable benefit
on conversation but only 50.2\% on toolagent, and that a learned ranker adds nothing on top of it on
either. So the recoverable benefit is not in per-item prediction. It is partly in \emph{when} the
decision is taken and partly --- on toolagent, worth 14 points --- in the order arrivals are admitted in,
which we do not currently model. What remains unclaimed on toolagent is large: half the reachable
benefit, and 81\% of the hindsight gap that arrival timing does not expose at all. Both have to come from
something the per-item formulation does not express: admission under contention, coordinated eviction
across concurrently-live sessions, or migration timed against idle windows. Identifying which of these
carries the gap, and building an evaluation that can see it, is the substantive next step, and it is a
scheduling question rather than a modelling one.

\paragraph{Calibration, and what a mis-specified cost model hides.}
Our tables are placeholders, and we saw concretely how misleading that is: under them the per-item optimal
action is identical for every reused block, which collapses a three-axis decision into a one-dimensional
sort and makes every policy look alike. Fitting the tables to a real vLLM+LMCache deployment --- the only
GPU-gated step in this work --- is a precondition for any claim about the representation and timing axes,
which our traces do not exercise at all.

\paragraph{Aggregate working set, not per-request footprint.}
Reactive sufficiency is governed by how much of the \emph{concurrently live} working set a tier holds, not
by one request's block count: holding tier size fixed in absolute blocks and raising the number of pooled
conversations still degrades the ratio. Characterizing that surface --- tier size against concurrency,
across sources --- would give an operator a sizing rule, which is the most directly actionable thing this
line of work could produce.

\paragraph{Composing placement with retrieval.}
\Cref{sec:placement-vs-retrieval} indicates the two are orthogonal in kind: placement sets latency and
cost, retrieval sets which facts reach the model. The open problem is the composition --- a retriever
running over a tiered store whose hot set a scheduler maintains, where a retrieval miss on a cold tier
costs latency rather than accuracy. Neither literature evaluates that joint object, and our own LoCoMo
measurement is far too small ($n=30$) to speak to it.

\paragraph{Sampling as a first-class experimental variable.}
Three of our ten defects were extrapolations from a single value of a parameter fixed for convenience, and
the most damaging reversed a headline claim's sign. For trace-driven work we would now require of
ourselves, and suggest for the area: randomized rather than prefix sampling, several disjoint draws with
intervals, and a reported sweep over every quantity the design fixes --- trace length, pool size, tier
size, split fraction --- before any claim that depends on one of them.

\paragraph{Safety and governance.}
A memory written by tools and other agents is an attack surface, and provenance must survive consolidation
so a fact can be traced, invalidated, or unlearned. A budget-constrained scheduler is a natural
enforcement point; doing that without giving back the frontier is unexplored.
